\section{Attaining the integral bound by handle cancellation}
\label{handle:appendix}\label{app:handles}

We prove the existence assertion in Theorem~\ref{morse:equality}.
Throughout, $W$ is a simply connected closed smooth $D$-manifold with
$D\geq9$.  The argument is the classical handle proof of Smale's theorem
\cite{Smale1963}, presented here with the low-index cancellation made
explicit, as in \cite[Appendix~A]{OpenAICompanion}.

We use the correspondence between Morse functions and ordered handle
decompositions, and the standard operations of handle slides, insertion
of cancelling pairs, and cancellation.  A handle of index $i$ is a copy
of $D^i\times D^{D-i}$ attached along
$S^{i-1}\times D^{D-i}$; its belt sphere in the new boundary is
$\{0\}\times S^{D-i-1}$.  A slide changes the integral basis of the
handles of a fixed index by an elementary addition.  Consecutive
handles of indices $i$ and $i+1$ cancel when the attaching sphere of the
latter meets the belt sphere of the former transversely in one point.
Attaching maps of subsequent handles are transported by these
operations, so the represented manifold is unchanged.
For the Morse correspondence, see \cite{Milnor1963}; for cancellation
and integral basis changes, see
\cite[Theorems~5.4 and~7.6]{Milnor1965}.

\subsection{Eliminating handles of index one and coindex one}

Start with an ordered handle decomposition.  Connectedness permits the
cancellation of all but one $0$-handle along a tree of $1$-handles.
The dual operation leaves one $D$-handle.  We now eliminate the remaining
$1$-handles, allowing insertion of cancelling $2/3$ pairs, as in
\cite[Theorem~8.1]{Milnor1965}.

Let $P$ be the union of the $0$- and $1$-handles.  It retracts onto a
graph, whose free generators correspond to the $1$-handles.  The map
$\pi_1(\partial P)\to\pi_1(P)$ is an isomorphism: the handles of $P$
relative to its boundary, read in reverse order, have indices $D-1$
and $D$.  Since $\pi_1(W)=0$ and handles of index at least three do
not change fundamental groups, the attaching words of the $2$-handles
normally generate this free group.

Fix a generator.  Insert a cancelling $2/3$ pair whose $2$-handle
initially attaches along a trivial circle in a small ball in $\partial P$.
Express the chosen generator as a product of conjugates of the old
attaching words and their inverses.  Sliding the new $2$-handle over
the old ones realizes this product: a slide is a band sum, the band
path gives the conjugation, and its orientation gives the choice of
word or inverse.  General position allows the bands to avoid the other
attaching circles; their tubular attaching regions may be shrunk.

The new attaching circle is now homotopic to a standard circle that
runs once over the chosen $1$-handle and avoids the others.  Choose the
standard circle disjoint from the other attaching circles.  In the
$(D-1)$-dimensional boundary, a generic homotopy between these circles
is an isotopy and avoids the other attaching circles: the expected
dimensions for self-intersections or intersections with another circle
in a one-parameter family are at most $3-(D-1)<0$.
Transport the framing along this isotopy.  The new $2$-handle now
cancels the chosen $1$-handle, since its attaching circle meets that
handle's belt sphere once.  Repetition removes all $1$-handles.

Apply the same construction to the reversed decomposition to eliminate
the $(D-1)$-handles.  Its inserted pairs have original indices
$D-3,D-2$, so no $1$-handles are reintroduced.  We have obtained one
handle in each of indices $0,D$, none in indices $1,D-1$, and all
remaining handles in indices $2,\ldots,D-2$.

\subsection{Locating an algebraically cancellable pair}

The manifold $W$ is orientable.  Choose orientations and let $C_*$ be
the integral handle complex.  Its boundary coefficients are algebraic
intersection numbers between attaching spheres and belt spheres in the
level between consecutive indices.  With
$B_i=\operatorname{im}\partial_{i+1}$, equation
\eqref{morse:chain-rank} gives
\begin{equation}\label{handle:excess}
 \sum_i\operatorname{rank}C_i-\lambda(W)
 =2\sum_i\bigl(\operatorname{rank}B_i-t_i(W)\bigr).
\end{equation}
Each summand on the right is nonnegative.  If the total number of
handles exceeds $\lambda(W)$, then
$\operatorname{rank}B_i>t_i(W)$ for some $2\leq i\leq D-3$;
the vanishing chain groups in degrees $1,D-1$ exclude other indices.

Because $C_i/Z_i\cong B_{i-1}$ is free, $Z_i=\ker\partial_i$ is
a direct summand of $C_i$.  Thus the nonzero Smith invariant factors
of $\partial_{i+1}$ are those of $B_i\subset Z_i$.  Exactly $t_i(W)$
of them are nonunits.  The strict inequality above supplies a unit
invariant factor.  Row and column operations, realized by handle slides,
relabeling, and orientation changes, therefore produce a boundary
coefficient equal to $1$.

These slides occur in connected levels.  Indeed, let $P_i$ be the union
of handles of index at most $i$, and put $L=\partial P_i$.
The space $P_i$ is simply connected because it has no $1$-handles.
Its reversed handles relative to $L$ have indices at least $D-i\geq3$,
so $L\to P_i$ induces isomorphisms on components and fundamental groups.
The same argument applies before the $i$-handles are attached.
Slide paths can avoid other attaching spheres, since those spheres have
codimension at least two; again their tubular regions may be shrunk.

We have found an attaching sphere $A\cong S^i$ and a belt sphere
$B\cong S^{D-1-i}$ in the simply connected level $L$, with algebraic
intersection number $1$.  We next arrange that they have exactly one
geometric intersection, which is the condition for handle cancellation.

\subsection{Removing opposite-sign intersections}

Suppose first that $i\geq3$ and $D-1-i\geq3$.
Join two opposite-sign intersections of $A$ and $B$ by an arc in each
sphere, avoiding the other intersection points.  The resulting Whitney
circle bounds a disk because $L$ is simply connected.
The high-dimensional Whitney trick removes the pair
\cite[Theorem~6.6]{Milnor1965}.
Both sphere dimensions are at least three, so general position keeps
the interior of the Whitney disk away from all attaching spheres of
the $(i+1)$-handles and all belt spheres of the $i$-handles.
The move preserves every other intersection at this level.
Repeating leaves just one intersection of $A$ and $B$.

When one sphere has dimension two, reverse the decomposition if
necessary so that $i=2$.  The level $L$ is then the boundary after
attaching the $2$-handles to the single $0$-handle.  Denote all their
belt spheres by $B_1,\ldots,B_h$, and set
\[
 V=L\setminus\bigcup_{j=1}^h B_j.
\]
We will remove a pair of intersections by moving a small disk in $A$
into $V$.  The special geometry of this complement makes that possible.

The boundary $L$ is obtained from $S^{D-1}$ by surgery on the attaching
circles.  Each surgery replaces $S^1\times D^{D-2}$ by
$D^2\times S^{D-3}$; removing its belt sphere leaves
$(D^2\setminus\{0\})\times S^{D-3}$.
This piece retracts onto its attaching boundary
$S^1\times S^{D-3}$.  Therefore $V$ has the homotopy type of the
exterior of the original attaching circles in $S^{D-1}$.
That exterior is simply connected: every loop bounds a disk in the
sphere, and general position makes the disk disjoint from the circles.
The complement of the circles retracts onto their exterior by radial
retraction in tubular neighborhoods.  We have proved
\begin{equation}\label{handle:complement-pi1}
 \pi_1(V)=\pi_1(L)=0.
\end{equation}

Choose two opposite-sign intersections of the attaching sphere
$A\cong S^2$ with one chosen belt sphere.  A thin disk
$\Delta\subset A$ surrounding an arc between them contains just these
two intersections and misses all other belt spheres.  Its boundary,
and an exterior collar in $A$, lie in $V$.
The oriented normal bundle of each belt sphere has rank two.
Excision and the Thom isomorphism give
\begin{equation}\label{handle:intersection-group}
 H_2(L,V;\Z)\cong\bigoplus_{j=1}^h\Z,
\end{equation}
with coordinates given by intersection numbers with the $B_j$.
All coordinates of $[\Delta]$ vanish, so $[\Delta]=0$.

This homological vanishing supplies a relative homotopy, without any
assumption on $\pi_2(L)$.  Indeed, \eqref{handle:complement-pi1} and
the exact sequences of the pair identify
\begin{align*}
 \pi_2(L,V)&\cong
 \operatorname{coker}\bigl(\pi_2(V)\longrightarrow\pi_2(L)\bigr),\\
 H_2(L,V;\Z)&\cong
 \operatorname{coker}\bigl(H_2(V;\Z)\longrightarrow H_2(L;\Z)\bigr).
\end{align*}
The degree-two Hurewicz isomorphisms for $V$ and $L$ identify these
cokernels.  Hence $\Delta$ can be homotoped into $V$ through maps whose
boundary stays in $V$.  Extend this boundary motion across the exterior
collar, fixing its outer edge.  We obtain a homotopy of $A$ supported
in the disk and collar; at the endpoint it removes the chosen pair and
preserves all other intersections with the belt spheres.

Finally this homotopy may be chosen to be an isotopy of the attaching
sphere.  Smooth it relative to the initial sphere and its stationary
part, and apply parametric general position.  For a one-parameter
family of maps of a $2$-sphere into the $(D-1)$-manifold $L$, the
expected dimension of self-intersections, or intersections with another
attaching $2$-sphere, is at most
\[
 2+2+1-(D-1)<0.
\]
The same general-position argument makes every map in the family an
immersion.  Since $D-1\geq8$, the relative version applies with ample
dimension to spare.  At the endpoint, the altered disk and collar form
a compact subset of the open set $V$, so a sufficiently small
approximation preserves their avoidance of all belt spheres.
Keep neighborhoods of the other transverse intersections stationary.
Intermediate intersections with belt spheres impose no restriction on
an isotopy of an attaching sphere.  Isotopy extension transports the
framing, while the other attaching tubes can be shrunk to remain
disjoint.  Repetition leaves exactly one intersection with the chosen
belt sphere.

\subsection{Finishing the count}

The chosen $i$-handle may now be ordered last among the $i$-handles,
and the chosen $(i+1)$-handle first among the $(i+1)$-handles.
They are consecutive and cancel.  Intersections with other handles do
not obstruct this operation; their attaching maps are carried along
by the cancellation diffeomorphism.
No handles of indices $1,D-1$ are introduced.

Each cancellation decreases the total count by two.  Whenever that
count exceeds $\lambda(W)$, equation~\eqref{handle:excess} supplies
another pair to which the preceding argument applies.  The process
therefore ends with exactly $\lambda(W)$ handles.  Equality in
\eqref{handle:excess} forces
$\operatorname{rank}B_i=t_i(W)$ for every $i$; substituting in
\eqref{morse:chain-rank} gives
\[
 \operatorname{rank}C_i=r_i(W)+t_i(W)+t_{i-1}(W).
\]
The Morse function associated with this decomposition has precisely
these index counts, proving Theorem~\ref{morse:equality}.
